\chapter{Introducción}
 
El procesamiento digital de imágenes constituye una de las áreas más demandantes en términos
computacionales dentro de la informática moderna. Si bien las bases matemáticas de muchas de
sus técnicas datan de siglos atrás, su aplicación práctica a escala masiva solo fue posible con
el avance tecnológico del hardware de las últimas décadas, que permitió realizar en tiempo
razonable la enorme cantidad de operaciones aritméticas involucradas.
 
Una imagen digital puede representarse como un arreglo bidimensional de píxeles. En el modelo
de color \textbf{RGB} (Red, Green, Blue) utilizado en este trabajo, cada píxel contiene tres
valores enteros en el rango $[0, 255]$, uno por canal de color. Esta representación permite
codificar hasta 16,8 millones de tonalidades distintas, superando la capacidad discriminatoria
del ojo humano. Los casos extremos son el negro $(0,0,0)$ y el blanco $(255,255,255)$.
 
El objetivo del presente trabajo es la \textbf{generación automática de imágenes tipo
\textit{cartoon}} a partir de una imagen original en formato RGB. Esta transformación busca
simplificar formas, resaltar contornos y reducir variaciones tonales, logrando una apariencia
similar a la de una ilustración o dibujo animado. Para ello se aplican dos procesos
independientes sobre la imagen original:
 
\begin{itemize}
  \item \textbf{Posterización:} reduce el rango total de niveles de color a un conjunto más
        pequeño de valores discretos. Cada píxel pasa a adoptar el valor del nivel más cercano
        dentro de ese rango reducido. En imágenes a color se posteriza cada canal por separado
        y luego se suman las imágenes resultantes.
 
  \item \textbf{Detección de bordes (borroneado):} pipeline de tres etapas secuencialmente
        dependientes entre sí:
        \begin{enumerate}
          \item \textit{Filtrado (blur):} suaviza la imagen para eliminar ruido.
          \item \textit{Resaltado (highlight):} amplifica las diferencias entre zonas de
                borde y no-borde mediante aproximación de gradiente (filtros Prewitt, Sobel
                u otros).
          \item \textit{Umbralado:} binariza el resultado del resaltado para obtener los
                bordes detectados.
        \end{enumerate}
\end{itemize}
 
La imagen cartoon final se obtiene sumando la salida de la posterización y la de la detección
de bordes.
 
Dado el volumen de cómputo involucrado, este trabajo explora y compara distintas \textbf{estrategias de
paralelización} con el fin de reducir el tiempo de procesamiento. Se implementan cuatro
versiones del algoritmo:
 
\begin{itemize}
  \item \textbf{Secuencial:} implementación de referencia, sin paralelismo.
  \item \textbf{Memoria compartida (OpenMP):} paralelismo a nivel de threads dentro de un
        único nodo mediante directivas.
  \item \textbf{Memoria distribuida (MPI):} paralelismo mediante pasaje de mensajes entre
        procesos.
  \item \textbf{Híbrido (MPI + OpenMP):} combinación de ambos paradigmas para explotar
        simultáneamente múltiples nodos y múltiples núcleos por nodo.
\end{itemize}
 
El análisis de desempeño se realiza sobre imágenes de tres tamaños
(800$\times$800, 2000$\times$2000 y 5000$\times$5000 píxeles), con filtros de convolución
de 3$\times$3 y 5$\times$5, y posterizado configurado con 3 y 9 rangos de color. Las métricas
utilizadas son el \textit{speedup} y la eficiencia paralela.